    \begin{minipage}{0.58\linewidth}
        Un enfant tire un chariot de $\qty{2.0}{\kg}$ sur une distance de
        $\qty{20}{\m}$, le sol faisant un angle $\beta$ de $\qty{10}{\degree}$
        par rapport à l'horizontale. La tension du fil à une direction faisant
        un angle $\alpha$ de $\qty{20}{\degree}$ par rapport au sol. Son
        intensité est de $\qty{9.2}{\N}$.

        On néglige les frottements.

    \end{minipage}
    \hfill
    \begin{minipage}{0.40\linewidth}
        \begin{figure}[H]
            \flushright
            \begin{tikzpicture}
                \def\a{20}
                \def\b{10}
                \def\F{1.2}
                \def\AC{6.3}
                \pgfmathparse{\AC*tan(\b)}\let\BC=\pgfmathresult
                \coordinate (A) at (0,0);
                \path[fill=black!10,draw,thick]
                    (A) -- ++(\AC,0) coordinate (C) -- ++(0,\BC)
                    coordinate (B) -- cycle;
                \node[solide={1cm/7mm},rotate=\b,anchor=south,
                    outer ysep=1.2pt] (S) at ($(A)!0.5!(B)$) {};
                \node[circle,fill,inner sep=0.8pt] (F) at (S.east) {};
                \draw[dashed] (F.center) -- ++(\b:2.8);
                \draw[thick,postaction={decorate},
                    decoration={markings,mark=at position 0.8 with
                    {\arrow{stealth}}}]
                    (F.center) -- ++(\b+\a:3);
                \draw[->,shorten >=0.4pt]
                    ([xshift=2cm]A) arc (0:\b:2cm)
                    node[pos=0.5,right,font=\footnotesize] {$\beta$};
                \draw[->,shorten >=0.4pt]
                    ([shift=(\b:1)]F.center) arc (\b:\b+\a:1)
                    node[pos=0.7,right,font=\footnotesize] {$\alpha$};
            \end{tikzpicture}
        \end{figure}
        
    \end{minipage}

\begin{enumerate}
    \item Déterminer les forces exercées sur le chariot.
    \item Calculer les travaux des forces extérieures exercées sur ce chariot.
\end{enumerate}

